%%
%% winter-lecture28.tex
%% 
%% Made by Alex Nelson <pqnelson@gmail.com>
%% Login   <alex@lisp>
%% 
%% Started on  2026-03-10T09:52:04-0700
%% Last update 2026-03-10T09:52:04-0700
%% 

\lecture{}

\begin{recall}
We defined the smash product $X\smashtimes Y=(X\times Y)/(X\wedgetimes Y)$.
We saw $\sphere{1}\smashtimes X\iso\Sigma X$.
\end{recall}

\begin{definition}
Given pointed compact Hausdorff spaces $(X,x_{0})$ and $(Y,y_{0})$, we
have the projections
\begin{subequations}
  \begin{align}
\operatorname{pr}_{X}\colon (X\times Y,x_{0}\times y_{0})\to(X,x_{0})\\
\intertext{and}
\operatorname{pr}_{Y}\colon (X\times Y,x_{0}\times y_{0})\to(Y,y_{0})
  \end{align}
\end{subequations}
which then induce a map called the \define{External Product}:
\begin{equation}
\begin{split}
\widetilde{K}(X)\otimes\widetilde{K}(Y)&\to\widetilde{K}(X\times Y)\\
a\otimes b&\mapsto\operatorname{pr}^{*}_{X}(a)\operatorname{pr}^{*}_{Y}(b),
\end{split}
\end{equation}
where we just us the usual multiplication in the K-theory on the
right-hand side. Note that
\begin{subequations}
  \begin{align}
\operatorname{pr}^{*}_{X}\colon\widetilde{K}(X)\to\widetilde{K}(X\times Y)\\
\intertext{and}
\operatorname{pr}^{*}_{Y}\colon\widetilde{K}(Y)\to\widetilde{K}(X\times Y)
  \end{align}
\end{subequations}
are the pullbacks of the projection operators, at the level of K-theory.
\end{definition}

\begin{remark}
I have seen the nLab claim the ``external product'' to be defined on
vector bundles $E_{1}\to X_{1}$ and $E_{2}\to X_{2}$ by
$E_{1}\boxtimes E_{2}\to X_{1}\times X_{2}$ to be the usual tensor
product of the two pullback bundles induced by pulling-back the projections
$\operatorname{pr}_{i}\colon X_{1}\times X_{2}\to X_{i}$.
\end{remark}

\begin{node}
Now, consider the long exact sequence obtained by
\begin{equation}
\begin{split}
X\smashtimes Y & \stackrel{i}{\into} X\times
Y\xrightarrow{q}X\wedgetimes Y\\
x & \mapsto x\times y_{0}\\
y & \mapsto x_{0}\times y
\end{split}
\end{equation}
We get a long exact sequence
\begin{equation}
\widetilde{K}^{1}(X\smashtimes Y)\to\widetilde{K}^{0}(X\wedgetimes
Y)\xrightarrow{q^{*}}\widetilde{K}^{0}(X\times
Y)\xrightarrow{i^{*}}\widetilde{K}^{0}(X\wedgetimes Y)\to0.
\end{equation}
We want to prove
\begin{equation}
\widetilde{K}^{0}(X\times Y)\iso\widetilde{K}^{0}(X\smashtimes
Y)\oplus\underbrace{\widetilde{K}^{0}(X\wedgetimes Y)}_{\iso\widetilde{K}(X)\oplus\widetilde{K}(Y)}
\end{equation}
Note that
\begin{equation}
\operatorname{pr}_{X}\circ i|_{X}\iso\id_{X},\quad\mbox{and}\quad
\operatorname{pr}_{X}\circ i|_{Y}\iso\id_{Y}.
\end{equation}
Then these combine to give the induced morphism
\begin{equation}
i^{*}\circ(\operatorname{pr}_{X}^{*}\oplus\operatorname{pr}_{Y}^{*})\colon
\widetilde{K}(X)\oplus\widetilde{K}(Y)\iso\widetilde{K}(X\wedgetimes Y)\to
\widetilde{K}(X\times Y)\to\widetilde{K}(X\smashtimes Y)
\end{equation}
And the claim is that this is the identity morphism.

Then
\begin{equation}
\operatorname{pr}_{X}^{*}\oplus\operatorname{pr}_{Y}^{*}\colon\widetilde{K}^{0}(X\wedgetimes Y)\to\widetilde{K}^{0}(X\times Y)
\end{equation}
is a section, so
\begin{equation}
i^{*}\circ(\operatorname{pr}_{X}^{*}\oplus\operatorname{pr}_{Y}^{*})=\id.
\end{equation}
We have
\begin{subequations}
  \begin{align}
(\Sigma i)^{*}\colon K^{1}(X\times Y)\to\widetilde{K}^{1}(X\wedgetimes Y)\\
    \intertext{and}
    \boundary\colon K^{1}(X\wedgetimes Y)\to K^{0}(X\smashtimes Y).
  \end{align}
\end{subequations}
We claim the connecting morphism $\boundary=0$ is just the zero map,
so we trim the long exact sequence into the short exact sequence
\begin{equation}
0\to\widetilde{K}^{0}(X\smashtimes
Y)\xrightarrow{q^{*}}\widetilde{K}^{0}(X\times
Y)\xrightarrow{i^{*}}\widetilde{K}(X\wedgetimes Y)\to0.
\end{equation}
This is because $i^{*}$ is surjective, so $\Sigma i^{*}$ is
surjective, so we must have $\boundary=0$.
\end{node}

\begin{node}
Under this decomposition,
\begin{equation}
\operatorname{pr}_{X}^{*}\operatorname{pr}_{Y}^{*}\colon\widetilde{K}(X)\oplus\widetilde{K}(Y)\to\widetilde{K}(X\times Y),
\end{equation}
and
\begin{equation}
K(X\times Y)\iso\underbrace{\widetilde{K}(X\smashtimes Y)}_{\sim\ker(i^{*})}\oplus\underbrace{\widetilde{K}(X)\oplus\widetilde{K}(Y)}_{\sim\Im(i^{*})}
\end{equation}
and $\operatorname{pr}_{X}^{*}\operatorname{pr}_{Y}^{*}$ lands in
$\widetilde{K}(X\smashtimes Y)$. That is to say,
\begin{equation}
i^{*}(\operatorname{pr}_{X}^{*}(-)\operatorname{pr}_{Y}^{*}(-))=0
\end{equation}
which is the same as saying
\begin{equation}
X\times\{y_{0}\}\into X\wedgetimes Y\stackrel{i}{\into}X\times Y\to0,
\end{equation}
so we get the external product
\begin{equation}
\boxtimes\colon\widetilde{K}(X)\otimes\widetilde{K}(Y)\to\widetilde{K}(X\smashtimes Y).
\end{equation}
For the usual K-theory $K(X)$ (= $\widetilde{K}(X)\oplus\ZZ$), we have
the multiplication
\begin{equation}
K(X)\otimes K(Y)\to K(X\times Y),
\end{equation}
which is a ring morphism.
\end{node}

\begin{theorem}[Fundamental product]
Let $X$ be a compact Hausdorff space. Then we have a ring morphism
\begin{equation}
\begin{split}
\mu\colon K(X)\times K(\sphere{2})&\to K(X\times\sphere{2})\\
a\otimes b&\mapsto(\operatorname{pr}_{1}^{*}a)(\operatorname{pr}_{2}^{*}b)
\end{split}
\end{equation}
is an isomorphism, where $K(\sphere{2})\iso\ZZ[H]/\bigl((H-1)^{2}\bigr)\iso\ZZ[\mathbf{1}]\oplus\ZZ[H]$.
\end{theorem}

\begin{claim}
Bott periodicity follows from the fundamental product theorem.
\end{claim}

\begin{proof}
We have
\begin{equation}
\mu\colon(\widetilde{K}(X)\oplus\ZZ)\otimes(\widetilde{K}(\sphere{2})\oplus\ZZ)\xrightarrow{\iso}\widetilde{K}(X\times\sphere{2})\oplus\ZZ.
\end{equation}
Then
\begin{equation}
\mu\colon(\widetilde{K}(X)\otimes\widetilde{K}(\sphere{2}))\oplus\underbrace{(\ZZ\otimes\widetilde{K}(\sphere{2}))}_{\iso\widetilde{K}(\sphere{2})}\oplus\underbrace{(\widetilde{K}(X)\otimes\ZZ)}_{\iso\widetilde{K}(X)}\oplus\ZZ\to\widetilde{K}(X\times\sphere{2})\oplus\ZZ,
\end{equation}
then $\mu$ induces an isomorphism between each component. For example,
$\widetilde{K}(\sphere{2})$ is sent to a trivial bundle over $X$ and
identity on itself, and so on.
\end{proof}

\begin{node}
We will follow a proof by Atiyah and Bott~\cite{atiyah1964periodicity},
which Hatcher's proof is close in spirit to. See also Husemoller's
\textit{Fiber Bundles} (third ed., Springer, Chapter 11).
\end{node}

\subsection{General Clutching Construction}

\begin{node}
Recall the usual clutching construction says vector bundles over
$\sphere{2}$ may be viwede as vector bundles over the hemispheres
(each of which is contractible), so all the important information
comes from the equator. The general version is as follows.
\end{node}

\begin{claim}
There is a bijective correspondence
\begin{equation}
\left\{\begin{matrix}\CC^{n}\mbox{-vector bundles}\\
\mbox{over }X\times\sphere{2}
\end{matrix}\right\}/\sim\quad\stackrel{1:1}{\longleftrightarrow}\quad\{(E,f)\}
\end{equation}
where $E$ is a $\CC^{n}$-bundle over $X$, with pullback bundle
\begin{equation}
\vcenter{\xymatrix{p^{*}E\ar[d]\ar[r] & E\ar[d]\\
X\times\sphere{1}\ar[r] & X}}
\end{equation}
and $f$ is a bundle isomorphism
\begin{equation}
\vcenter{\xymatrix{p^{*}E\ar[rr]^{f}\ar[dr]& & p^{*}E\ar[dl]\\
& X\times\sphere{1} &}}.
\end{equation}
We recover the usual clutching construction when $X=\{x\}$ is a point.
\end{claim}

The proof is roughly the same as before.

\begin{proof}
Given a pair $(E,f)$ we will produce an isomorphism class
\begin{equation}
[E,f]\mapsfrom(E,f)
\end{equation}
as follows:
$[E,f]$ is found by gluing two copies of $E\times\disk{2}_{\pm}\to X\times\disk{2}_{\pm}$
along the ``equator'' $E\times\sphere{1}$ using the automorphism $f$
\begin{equation}
(E\times\disk{2}_{+})\sqcup(E\times\disk{2}_{-})/(e,z)\sim f(e,z)
\end{equation}
where $e\in E$ and $z\in\sphere{1}$. The isomorphism class $[E,f]$
depends only on the homotopy class of $f$ among bundle automorphisms
of $E\times\sphere{1}$ (this is homework).

Now, to prove this is a bijection, we can construct the inverse map as
follows: Given a $\CC^{n}$-bundle $\pi\colon E'\to X\times\sphere{2}$
over $X\times\sphere{2}=(X\times\disk{2}_{+})\cup(X\times\disk{2}_{-})$,
we get $E'_{+}\to X\times\disk{2}_{+}$ and $E'_{-}\to X\times\disk{2}_{-}$
by restriction of $\pi$. Let
\begin{equation}
\begin{split}
i\colon &X\into X\times\disk{2}_{+}\\
&x\mapsto(x,p)
\end{split}
\end{equation}
for some fixed $p\in\disk{2}_{+}$, and
\begin{equation}
\begin{split}
r\colon&X\times\disk{2}_{+}\to X\\
&(x,d)\mapsto x
\end{split}
\end{equation}
for the retraction. Then
\begin{equation}
i\circ r\homotopic \id_{X\times\disk{2}_{+}},
\end{equation}
so we get
\begin{equation}
E'^{+}\iso(i\circ r)^{*}E'^{+}\iso (r^{*}\circ i^{*})E'^{+}
\end{equation}
We call $i^{*}E'$ a $\CC^{n}$-bundle $E\to X$, then
\begin{equation}
E'^{+}\iso(E\times\disk{2}_{+}\to X\times\disk{2}_{+}).
\end{equation}
We have a bundle isomorphism $\varphi_{+}$
\begin{equation}
\vcenter{\xymatrix{E'_{+}\ar[dr]\ar[rr]^{\varphi_{+}} & & E\times\disk{2}_{+}\ar[dl]\\
& X\times\disk{2}_{+} &}}
\end{equation}
and similarly, we can construct
\begin{equation}
\vcenter{\xymatrix{E'_{-}\ar[dr]\ar[rr]^{\varphi_{-}} & & E\times\disk{2}_{-}\ar[dl]\\
& X\times\disk{2}_{-} &}}
\end{equation}
Then we get the clutching map $f$ by composing the two trivializations
over $X\times\sphere{1}$
\begin{equation}
\vcenter{\xymatrix{E\times\sphere{1}\ar[r]^{(\varphi_{-})^{-1}}\ar[dr]
    & E'|_{\sphere{1}} \ar[d] \ar[r]^{\varphi_{+}} & E\times\sphere{1}\ar[dl]\\
& X\times\sphere{1} &}}
\end{equation}
where $f=\varphi_{+}\circ\varphi_{-}^{-1}$.
\end{proof}

\begin{remark}
We can view the clutching map $f\colon E\times\sphere{1}\to E\times\sphere{1}$
as an $\sphere{1}$-parametrized family of $\CC^{n}$-bundle
automorphisms $E\to E$ so we may write $f$ as
\begin{equation}
f_{z}\colon E\to E
\end{equation}
for $z\in\sphere{1}$.
\end{remark}

\begin{node}
\textbf{Outline of proof of fundamental product theorem}.
\begin{enumerate}
\item Up to homotopy, we cna assume $f_{z}$ takes the form of a
  Laurent polynomial in $z$,
\begin{equation}
f_{z}(x)=\sum_{|k|<n}a_{k}(x)z^{k},
\end{equation}
by Fourier analysis.
\item (In the K-theory world) We can assume $f_{z}$ is a polynomial in
  $z$ by shifting the negative powers of $z$ away
\item Using $K(\sphere{2})\iso\ZZ[H]/(H-1)^{2}$ we can reduce to the
  case where $f_{z}$ is linear
\begin{equation}
f_{z}(x)=a_{0}(x)+a_{1}(x)z=a(x)+b(x)z.
\end{equation}
\item Given $[E,a(x)+zb(x)]$, this is the same (up to trivialization)
  to some splitting
\begin{subequations}
\begin{equation}
E_{+}\oplus E_{-}\iso E
\end{equation}
and
\begin{equation}
[E,a(x)+zb(x)]\iso\underbrace{[E_{+},a(x)]}_{\sim K(X)\otimes\ZZ[\mathbf{1}]}+\underbrace{[E_{-},zb(x)]}_{\sim K(X)\otimes\ZZ[H]}
\end{equation}
\end{subequations}
\end{enumerate}
\end{node}